% eksperymentalne tłumaczenie na włoski

%

\subsection{Aksjomaty leżenia pomiędzy} % lang-pl
\subsection{Assiomi di ordinamento} % lang-it

\index{aksjomat!leżenia pomiędzy} % lang-pl
\index{assioma!di ordinamento} % lang-it

O relacji między aksjomatami B1-B4 i uporządkowaniem ciała, nad którym pracujemy, rozwodzi się Hartshorne \cite[s. 135-140]{hartshorne_2000} % lang-pl
Hartshorne \cite[pp.~135--140]{hartshorne_2000} discute dettagliatamente la relazione tra gli assiomi B1--B4 e l'ordinamento del campo su cui si lavora. % lang-it

\begin{axiom}
[leżenia pomiędzy, B1] % lang-pl
[di ordinamento, B1] % lang-it
    Jeśli punkt $B$ leży między punktami $A$ i $C$, to punkty $A$, $B$, $C$ są różnymi punktami tej samej prostej oraz punkt $B$ leży także między punktami $C$ i $A$. % lang-pl
    Se il punto $B$ si trova tra i punti $A$ e $C$, allora $A$, $B$ e $C$ sono punti distinti della stessa retta e inoltre $B$ si trova anche tra $C$ e $A$. % lang-it
\end{axiom}

\begin{axiom}
[leżenia pomiędzy, B2] % lang-pl
[di ordinamento, B2] % lang-it
    Dla każdej pary punktów $A$ i $B$ istnieje punkt $C$ taki, że punkt $B$ leży między punktami $A$ i $C$. % lang-pl
    Per ogni coppia di punti $A$ e $B$ esiste un punto $C$ tale che $B$ si trovi tra $A$ e $C$. % lang-it
\end{axiom}

\begin{axiom}
[leżenia pomiędzy, B3] % lang-pl
[di ordinamento, B3] % lang-it
    Spośród trzech punktów leżących na prostej, dokładnie jeden leży pomiędzy pozostałymi dwoma. % lang-pl
    Fra tre punti allineati, esattamente uno si trova tra gli altri due. % lang-it
\end{axiom}

\begin{axiom}
[leżenia pomiędzy, B4] % lang-pl
[di ordinamento, B4] % lang-it
    Niech $A$, $B$ i $C$ będą trzema niewspółliniowymi punktami, zaś $l$ prostą, która nie przechodzi przez żaden z nich. % lang-pl
    Siano $A$, $B$ e $C$ tre punti non allineati e sia $l$ una retta che non passa per nessuno di essi. % lang-it

    Jeśli prosta $l$ przechodzi przez punkt między punktami $A$ i $B$, to przechodzi też przez punkt między punktami $A$ i $C$ albo $B$ i $C$, ale nie przez obydwa. % lang-pl
    Se la retta $l$ passa per un punto compreso tra $A$ e $B$, allora passa anche per un punto compreso tra $A$ e $C$ oppure tra $B$ e $C$, ma non per entrambi. % lang-it
\end{axiom}

Powyższy aksjomat nazywany jest też aksjomatem Pascha, ponieważ Moritz Pasch \cite{pasch_1882} przyłapał dopiero w 1882 roku geometrów całego świata na tym, że korzystali z takiej przesłanki. % lang-pl
L'assioma precedente è chiamato anche assioma di Pasch, poiché soltanto nel 1882 Moritz Pasch \cite{pasch_1882} fece notare che i geometri di tutto il mondo stavano già utilizzando implicitamente un'ipotesi di questo tipo. % lang-it

\index{aksjomat!Pascha} % lang-pl
\index{assioma!di Pasch} % lang-it

\index[persons]{Pasch, Moritz}%

Aksjomatu Pascha (B4) nie wolno mylić z aksjomatem Veblena-Younga geometrii rzutowej\footnote{Jeśli prosta przecina dwa boki trójkąta, to przecina także jego trzeci bok. \index{aksjomat!Veblena-Younga}}! % lang-pl
L'assioma di Pasch (B4) non va confuso con l'assioma di Veblen--Young della geometria proiettiva\footnote{Se una retta interseca due lati di un triangolo, allora interseca anche il terzo lato. \index{assioma!di Veblen--Young}}! % lang-it

% PRZEJRZANO: https://en.wikipedia.org/wiki/Pasch%27s_axiom

\begin{proposition}
    Z aksojmatów I1 do I3, B1 do B4 wynika, że każda prosta ma nieskończenie wiele punktów. % lang-pl
    Dagli assiomi I1--I3 e B1--B4 segue che ogni retta contiene infiniti punti. % lang-it
\end{proposition}

\begin{definition}
[odcinek] % lang-pl
[segmento] % lang-it
    Niech $A$, $B$ będą punktami. % lang-pl
    Siano $A$ e $B$ due punti. % lang-it

    Zbiór złożony z punktów $A$, $B$ oraz punktów, które leżą między nimi, nazywamy odcinkiem i oznaczamy: % lang-pl
    L'insieme formato dai punti $A$, $B$ e da tutti i punti che si trovano tra essi si chiama segmento e si indica con: % lang-it

\index{odcinek} % lang-pl
\index{segmento} % lang-it

    \begin{equation}
        \overline{AB}.
    \end{equation}
\end{definition}

Jeżeli nie prowadzi to do nieporozumień, kreskę pomijamy. % lang-pl
Quando non vi è rischio di ambiguità, la barra viene omessa. % lang-it

\begin{definition}
[trójkąt] % lang-pl
[triangolo] % lang-it
    Niech $A$, $B$, $C$ będą punktami. % lang-pl
    Siano $A$, $B$ e $C$ tre punti. % lang-it

    Sumę odcinków $AB$, $BC$, $AC$ nazywamy trójkątem, wspomniane odcinki -- jego bokami, zaś punkty $A$, $B$ i $C$ -- wierzchołkami. % lang-pl
    L'unione dei segmenti $AB$, $BC$ e $AC$ si chiama triangolo; tali segmenti sono i suoi lati e i punti $A$, $B$ e $C$ i suoi vertici. % lang-it

\index{trójkąt} % lang-pl
\index{triangolo} % lang-it
\end{definition} % Hartshorne 74

Trójkąty zazwyczaj oznaczamy tak: $\triangle ABC$ i nie odróżniamy go od trójkąta z wnętrzem. % lang-pl
Di solito indichiamo un triangolo con $\triangle ABC$ e non distinguiamo il bordo del triangolo dalla figura comprensiva del suo interno. % lang-it

Mamy podział na % lang-pl
Si distinguono % lang-it
\begin{itemize}
    \item trójkąty ostrokątne (z trzema kątami ostrymi), % lang-pl
    \item prostokątne (z jednym kątem prostym), % lang-pl
    \item rozwartokątne (z jednym kątem większym niż prosty), % lang-pl
    \item triangoli acutangoli (con tre angoli acuti), % lang-it
    \item rettangoli (con un angolo retto), % lang-it
    \item ottusangoli (con un angolo ottuso), % lang-it
\end{itemize}

ale też na: % lang-pl
e inoltre % lang-it

\begin{itemize}
\item równoboczne (z trzema bokami, które są przystające), % lang-pl
\item równoramienne (z dwoma przystającymi bokami zwanymi ramionami, wtedy trzeci bok nazywa się podstawą) albo % lang-pl
\item różnoboczne, z angielskiego \emph{scalene}. % lang-pl
\item equilateri (con tre lati congruenti), % lang-it
\item isosceli (con due lati congruenti detti lati obliqui; il terzo lato è chiamato base) oppure % lang-it
\item scaleni. % lang-it
\end{itemize}

\begin{proposition}
    Niech $l$ będzie prostą. % lang-pl
    Sia $l$ una retta. % lang-it

    Wtedy zbiór punktów, które nie leżą na prostej $l$ można rozbić na dwa niepuste zbiory $S_1$, $S_2$ takie, że: dwa punkty, które nie leżą na prostej $l$, należą do tego samego zbioru ($S_1$ lub $S_2$) wtedy i tylko wtedy, gdy odcinek $AB$ nie przecina prostej $l$. % lang-pl
    Allora l'insieme dei punti che non appartengono alla retta $l$ può essere suddiviso in due insiemi non vuoti $S_1$, $S_2$ tali che due punti esterni a $l$ appartengano allo stesso insieme ($S_1$ oppure $S_2$) se e solo se il segmento $AB$ non interseca la retta $l$. % lang-it
\end{proposition} % Hartshorne 74

Zbiory $S_1$, $S_2$ nazywamy stronami prostej $l$. % lang-pl
Gli insiemi $S_1$ e $S_2$ sono detti semipiani determinati dalla retta $l$. % lang-it

\index{prosta!dwie strony prostej} % lang-pl
\index{retta!due lati di una retta} % lang-it

Podobnie punkt wyznacza na prostej dwa zbiory, które leżą po różnych stronach tego punktu. % lang-pl
Analogamente, un punto determina sulla retta due insiemi situati da parti opposte rispetto a quel punto. % lang-it

\begin{definition}
[półprosta] % lang-pl
[semiretta] % lang-it
    Niech $A$, $B$ będą punktami. % lang-pl
    Siano $A$ e $B$ due punti. % lang-it

    Zbiór złożony z punktów $A$, $B$ oraz punktów, które leżą po tej samej stronie punktu $A$ na prostej $AB$ co punkt $B$, nazywamy półprostą nie ma sensownego oznaczenia % lang-pl
    L'insieme formato dai punti $A$, $B$ e da tutti i punti che si trovano sulla retta $AB$ dalla stessa parte di $A$ in cui si trova $B$ è chiamato semiretta; non esiste una notazione universalmente soddisfacente % lang-it

\index{półprosta} % lang-pl
\index{semiretta} % lang-it

    % \begin{equation}
        % \overrightarrow{AB}. % AB^{\rightarrow} ?
    % \end{equation}
\end{definition} % Hartshorne 77

% Ponownie, strzałkę psującą interlinię często pomijamy.

\begin{definition}
[kąt] % lang-pl
[angolo] % lang-it
    Sumę dwóch półprostych $AB$, $AC$, które nie leżą na jednej prostej, nazywamy kątem, zaś punkt $A$ wierzchołkiem tego kąta. % lang-pl
    L'unione di due semirette $AB$ e $AC$ che non giacciono sulla stessa retta si chiama angolo; il punto $A$ è il vertice dell'angolo. % lang-it

    Wnętrze kąta $\angle BACS$ składa się z tych punktów $D$ takich, że $D$ i $C$ leżą po tej samej stronie prostej $AB$ oraz $D$ i $B$ leżą po tej samej stronie prostej $AC$. % lang-pl
    L'interno dell'angolo $\angle BACS$ è costituito dai punti $D$ tali che $D$ e $C$ si trovino dalla stessa parte della retta $AB$ e che $D$ e $B$ si trovino dalla stessa parte della retta $AC$. % lang-it

    \index{kąt} % lang-pl
    \index{angolo} % lang-it
\end{definition} % Hartshorne 77

W myśl tej definicji, nie ma kąta zerowego, wklęsłego ani półpełnego. % lang-pl
Secondo questa definizione non esistono angoli nulli, concavi né piatti. % lang-it

\index{kąt!zerowy} % lang-pl
\index{angolo!nullo} % lang-it

\index{kąt!półpełny} % lang-pl
\index{angolo!piatto} % lang-it

\index{kąt!wklęsły} % lang-pl
\index{angolo!concavo} % lang-it

Wnętrze trójkąta $ABC$ to część wspólna wnętrz kątów $\angle ABC$, $\angle BCA$, $\angle CAB$; jest wypukłe (jeśli dwa punkty końce odcinka należą do wnętrza, to również każdy inny punkt tego odcinka znajduje się tam) i niepuste. % lang-pl
L'interno del triangolo $ABC$ è l'intersezione degli interni degli angoli $\angle ABC$, $\angle BCA$ e $\angle CAB$; esso è convesso (se gli estremi di un segmento appartengono all'interno, allora vi appartiene anche ogni punto del segmento) e non vuoto. % lang-it

%